\chapter[Probabilistic Setting]{Probabilistic Setting and Universal Optimality}\label{ch:sec:background}

To study encryption systems geometrically we first need a precise
probabilistic model of what an encryption system is.  What are the
ingredients?  Messages, keys, ciphertexts---each with a probability
distribution---and a deterministic rule that, given a message and a key,
produces a ciphertext.  This chapter formalises these ingredients and
introduces the central notion of \emph{universal optimality}: the
requirement that, after observing the ciphertext, every compatible
message be equally likely.

The probabilistic model of this chapter underpins every subsequent construction: the WIS factorisation, the gauge quotient, the posterior manifold, and the UOD all depend on the prior $\pi$, the message space $\Mset$, and the ciphertext space $\Cset$ introduced here.


\section{Finite probability spaces and notation}\label{ch:sec:finite-probability-spaces-and-notation}

Let $\Mset$, $\Kset$, and $\Cset$ be nonempty finite sets, called the
message, key, and ciphertext spaces, respectively.  Let
$\pi:\Mset\to(0,1]$, $\sum_m\pi(m)=1$, be a prior of full support.
The key is uniform on $\Kset$ and independent of the message:
\[
\Pr[K=k]=\frac1{|\Kset|},\qquad
\Pr[M=m,K=k]=\frac{\pi(m)}{|\Kset|}.
\]

\subsection{Deterministic encoders and the induced channel}\label{ch:sec:deterministic-encoders-and-the-induced-channel}

\begin{definition}\label{ch:def:encoder-basics}
A \emph{deterministic Shannon encoder} is a map
$E:\Mset\times\Kset\longrightarrow\Cset.$
The ciphertext random variable is $C=E(M,K)$.
\end{definition}

For $m\in\Mset$ and $c\in\Cset$,
\[
\Pr[C=c\mid M=m]
=\frac{|\{k\in\Kset:E(m,k)=c\}|}{|\Kset|}.
\]
The numerator will be formalised as the multiplicity matrix
$n_{mc}=|\{k:E(m,k)=c\}|$ in Section~\ref{ch:sec:wis}.

\section{Support, incidence, and cloaks}\label{ch:sec:support-incidence-and-cloaks}

\begin{definition}\label{ch:def:cloak-basics}
For $c\in\Cset$, the \emph{cloak} of $c$ is
$\cloak(c)=\{m\in\Mset:\Pr[C=c\mid M=m]>0\}.$
We say $m$ and $c$ are \emph{incident} ($m\sim c$) when $m\in\cloak(c)$.
\end{definition}

A ciphertext is \emph{active} when $\Pr[C=c]>0$ (its cloak is nonempty).

\subsection{Bayes' rule}\label{ch:sec:bayes-rule}

For every active ciphertext $c$ and every message $m$,
\begin{equation}\label{ch:eq:bayes-multiplicity-preview}
\Pr[M=m\mid C=c]
=\frac{\pi(m)}{|\Kset|\,\Pr[C=c]}\,
|\{k\in\Kset:E(m,k)=c\}|.
\end{equation}
Equation~\eqref{ch:eq:bayes-multiplicity-preview} is the bridge from
posterior symmetry to a combinatorial factorisation.

\section{Universal optimality relative to a prior}\label{ch:sec:universal-optimality-relative-to-a-prior}

Shannon's perfect secrecy~\cite{shannon1949} requires the posterior to
coincide with the prior for every ciphertext.  The present work studies
a weaker---but still highly structured---posterior symmetry: conditional
on the observed ciphertext, all messages in its cloak are equiprobable.
Following Biondi, Given-Wilson and
Legay~\cite{biondi2018apollonian}, we refer to this property as
\emph{universal optimality}.

\begin{definition}\label{ch:def:universal-optimality-basics}
A deterministic Shannon encoder is \emph{universally optimal relative
to $\pi$} if, for every active ciphertext $c$,
\[
\Pr[M=m\mid C=c]=\frac1{|\cloak(c)|}
\qquad\text{for every }m\in\cloak(c).
\]
\end{definition}

This is a posterior condition: after observing $c$, all compatible
messages are equiprobable.  It does not assert that the prior itself is
uniform.

\begin{remark}
Universal optimality is weaker than Shannon perfect secrecy: the former
requires $\Pr[M=m\mid C=c]$ to be uniform over the cloak, whereas the
latter requires $\Pr[M=m\mid C=c]=\pi(m)$ for every message.  The two
notions coincide only when the prior is uniform on every cloak and
every message is incident with every ciphertext (complete incidence):
otherwise a message outside the cloak of some ciphertext has zero
posterior probability there while its prior may be positive.
\end{remark}

A simple example is the one-time pad with uniformly distributed keys:
every ciphertext is compatible with all plaintexts, and the posterior
remains uniform.  Example~1 in Section~\ref{ch:sec:worked-examples}
provides a finite instance.

\section{First consequence of posterior uniformity}\label{ch:sec:first-consequence-of-posterior-uniformity}

Combining Definition~\ref{ch:def:universal-optimality-basics} with
Bayes' rule shows that, for each active $c$, the quantity
$\pi(m)\Pr[C=c\mid M=m]$ is independent of $m\in\cloak(c)$.
Equivalently, $\pi(m)n_{mc}$ is constant along every cloak.
Section~\ref{ch:sec:wis} records this in matrix form, and
Chapter~\ref{ch:sec:representation} turns it into the representation
theorem.

\paragraph{Running example: AES.}
Our main running example, the AES S-box (Section~\ref{ch:sec:aes-and-perfect-secrecy}), illustrates this simplest case.  The message and ciphertext spaces coincide ($|\Mset|=|\Cset|=2^{128}$), the incidence is complete ($I=\Mset\times\Cset$), and every ciphertext has a single pre-image: $n_{mc}=1$ for all $(m,c)$.  By Bayes' rule the posterior equals the prior, $\Pr[M=m\mid C=c]=\pi(m)$, on the full message space; with a uniform prior it is uniform on the cloak, which is all of $\mathcal M$, and universal optimality holds with $\mathcal D\equiv0$.

\paragraph{Scope.}
This book is restricted to finite spaces and deterministic encoders.
Probabilistic encoders and infinite spaces are left for future work.
\endinput